\chapter{Splint Application}

\section*{What Is This Test or Procedure?}
Splint application supports and immobilizes an injured body part while allowing room for swelling. It is commonly used for fractures, sprains, dislocations, tendon injuries, and painful joints.

\section*{How It Is Performed}
The clinician checks the skin, circulation, sensation, and movement before applying the splint. Padding and a rigid or flexible support are placed along the injured area and secured without excessive tightness. The limb is rechecked afterward.

\section*{Common Types}
Examples include short-arm, sugar-tong, thumb-spica, ulnar-gutter, volar, posterior ankle, and finger splints. The choice depends on the injury and the need to control movement while preserving swelling room.

\section*{Preparation and Aftercare}
Remove rings and tight jewelry when swelling is possible. Keep the splint dry and elevated as instructed. Do not alter the splint without medical advice. Move any joints that are not immobilized when permitted.

\section*{Risks and Warning Signs}
Pressure injury, skin irritation, stiffness, nerve compression, and reduced circulation can occur if a splint is too tight or poorly fitted. Increasing pain, numbness, tingling, pale or blue fingers or toes, severe swelling, or inability to move the digits requires prompt medical assessment.

\section*{Follow-up}
Follow-up may include repeat examination, imaging when needed, adjustment or replacement of the splint, and transition to a cast, brace, or rehabilitation program as healing progresses.

\section*{References}
\begin{enumerate}
  \item MedlinePlus. Splints. U.S. National Library of Medicine; 2025.
  \item Current Medical Diagnosis and Treatment. McGraw-Hill; 2025.
\end{enumerate}
\clearpage
